%% THERMODYNAMIQUE %%
%
\section{Masse volumique de l'air sec}

\itemstyle{ Op\'erateur~: } {\sf DensAirsec}
\index{DensAirsec}

%%%%%OBJET%%%%%
\itemstyle{ Objet~: }
Calcul de la masse volumique (�densit\'e�) de l'air sec, \`a 
partir de la temp\'erature et de la pression.

%%%%%ARGUMENTS%%%%%
\bigskip
%%%%%ENTREES%%%%
\begin{tabarg}
\multicolumn{3}{l}{\titretab{ Entr\'ees}}\\
\hline
$P_s$ &  Vecteur, {\tt float} &Pression ($hPa$) \\
$T_s$ & Vecteur, {\tt float} & Temp\'erature de l'air ($K$ ou $^\circ C$)\\
\end{tabarg}

\medskip
%%%%%SORTIES%%%%
\begin{tabarg}
\multicolumn{2}{l}{\titretab{ Sorties}}\\
\hline
$Dens$ & Vecteur, {\tt float} & Densit\'e de l'air sec ($kg/m^3$)\\
\end{tabarg}


%%%%%FORMULE%%%%
\itemstyle{ Formule }

%% calcul de la densite de l'air sec
$T_s$ est converti en Kelvin, si besoin est, puis :
\[
Dens = \frac{P_s\times 100}{R_a \times T_s}
\]
avec $R_a = 287.05 J.kg^{-1}.K^{-1} $.


%%%%%Commentaires%%%%%
\itemstyle{ Commentaires~:}

Si on passe une temp�rature virtuelle en argument (� la place de T),
on obtient (par d�finition m�me de Tv) la masse volumique de
l'air humide.

%%%%%AUTEUR%%%%%
\itemstyle{ Auteur~:}
CNRM/GMEI/TRAMM

%%%%%REFERENCES%%%%%
\itemstyle{ R\'ef\'erences~:}
Equation d'�tat d'un gaz parfait, Triplet-Roche \cite{Triplet-Roche},
page 34.
